\documentclass[12pt,a4paper]{report}
\usepackage[utf8]{inputenc}
\usepackage[spanish]{babel}
\usepackage{geometry}
\geometry{a4paper, margin=2.5cm}
\usepackage{graphicx}
\usepackage{hyperref}
\usepackage{titlesec}
\usepackage{booktabs}
\usepackage{amsmath}
\usepackage{fancyhdr}

\hypersetup{
    colorlinks=true,
    linkcolor=blue,
    filecolor=magenta,      
    urlcolor=cyan,
    pdftitle={Memoria Técnica - Hospitales},
    pdfpagemode=FullScreen,
}

\pagestyle{fancy}
\fancyhf{}
\rhead{Memoria Técnica}
\lhead{Proyecto Hospitales}
\cfoot{\thepage}

\title{
    \vspace{2cm}
    \Huge \textbf{Memoria Técnica de Proyecto} \\
    \vspace{1cm}
    \large Gestión y Optimización en Hospitales \\
    \vspace{2cm}
}
\author{\textbf{Autor:} frealr \& Antigravity}
\date{\today}

\begin{document}

\maketitle

\tableofcontents
\newpage

\chapter{Introducción}
\section{Contexto del Proyecto}
Este documento constituye la memoria técnica del proyecto de gestión y optimización de recursos en hospitales. Se describen los objetivos, el diseño del sistema y los resultados preliminares.

\section{Objetivos}
\begin{itemize}
    \item Optimizar la distribución de recursos hospitalarios.
    \item Mejorar la eficiencia en el flujo de pacientes.
    \item Diseñar un sistema de monitoreo en tiempo real.
\end{itemize}

\chapter{Descripción Técnica}
\section{Arquitectura del Sistema}
El sistema propuesto consta de varios módulos interconectados, diseñados para procesar datos en tiempo real.

\section{Especificaciones de Hardware y Software}
A continuación se detallan los requisitos técnicos del sistema:
\begin{table}[h]
\centering
\begin{tabular}{lll}
\toprule
\textbf{Componente} & \textbf{Tecnología} & \textbf{Función} \\
\midrule
Base de Datos & PostgreSQL & Almacenamiento persistente \\
Servidor Backend & Node.js / Python & Lógica de negocio y APIs \\
Frontend & React / Next.js & Panel de control de usuarios \\
\bottomrule
\end{tabular}
\caption{Requisitos tecnológicos de la solución.}
\label{tab:tecnologias}
\end{table}

\chapter{Conclusiones y Trabajo Futuro}
Se ha establecido la base del proyecto y se espera continuar con el desarrollo de los módulos de inteligencia artificial para la predicción de demanda de camas.

\end{document}
